\documentclass[10pt,a4paper]{amsart}
\usepackage[margin=19mm]{geometry}
\usepackage{amsmath,amssymb,mathtools}
\usepackage[scr=rsfs,cal=euler]{mathalfa}
\usepackage{xfrac}
\usepackage[unicode,hidelinks,bookmarks=false]{hyperref}
\newcommand{\ie}{i.e.\ }
\newcommand{\cf}{cf.\ }
\newcommand{\resp}{respectively\ }
\newcommand{\Subsection}{\subsection}
\providecommand{\nobreakdash}{\nobreak\hbox{-}}
\setlength{\emergencystretch}{2em}
\multlinegap0pt
\usepackage{bm}
\usepackage{bbm}
\usepackage{dsfont}
\usepackage{xspace}
\usepackage{cancel}
\usepackage{mathtools}
\usepackage{amsthm}
\makeatletter
\@ifundefined{lemma}{\newtheorem{lemma}{Lemma}}{}
\makeatother
\def\phi{\varphi}
\title[Rethinking Multi-User Communication in Semantic Domain: Enha]{Rethinking Multi-User Communication in Semantic Domain: Enhanced OMDMA by Shuffle-Based Orthogonalization and Diffusion Denoising}
\author{Maojun Zhang, Guangxu Zhu, Xiaoming Chen, Kaibin Huang, Zhaoyang Zhang}
\date{Source: arXiv:2507.20477v1 (CC BY 4.0)}
\begin{document}
\maketitle

\noindent\emph{Source and licence.} Rethinking Multi-User Communication in Semantic Domain: Enhanced OMDMA by Shuffle-Based Orthogonalization and Diffusion Denoising. Version arXiv:2507.20477v1 (CC BY 4.0), distributed under the Creative Commons Attribution 4.0 International licence, as recorded in the arXiv licence metadata for this exact version and shown on the abstract page https://arxiv.org/abs/2507.20477v1. Full rights notice: licenses/arxiv-2507.20477-CC-BY-4.0.txt.\par\smallskip

\noindent\textit{Editorial scope.} Counted unit: the Lemma whose source label is \texttt{lemma: received latent feature} in the article (source0.tex lines 555--563, source label \texttt{lemma: received latent feature}), delivered together with the article's own complete original proof of it (source0.tex lines 564--592, one \texttt{proof} environment of 29 lines). No mathematics is written, repaired or re-derived: statement and proof are the article's own text line for line. Required context retained uncounted: none. Every definition, display and label of those lines is retained unchanged. Omitted from this package and not redistributed: 1--554, 593--1217. Nothing inside a retained excerpt is omitted or altered: every retained body is delivered exactly as the article's lines are written, so no omission receipt of any kind accompanies this package. Citations: the retained excerpts carry no citation command at all, so no bibliography entry of the article is transcribed and no reference is redistributed. Reference adaptation: none was needed and none was made; every reference inside the delivered unit resolves inside it, so no unresolved reference or citation is produced. Printed slips kept literal: no printed word, symbol, hypothesis or number of the counted statement or of its proof was altered. Layout and class: the article's own class is \texttt{IEEEtran} with packages including fancyhdr, epsfig, threeparttable, epsf, amsthm, amsmath, amssymb, amsfonts, cite, dsfont, subfigure, color, soul, accents; the wrapper is the lane's pinned \texttt{amsart} editorial wrapper at 10pt on A4, which restates the theorem-like environments the retained text needs and adds the article's own macro definitions that the retained text needs (\texttt{\textbackslash phi}), with extra wrapper packages (bm, bbm, dsfont, xspace, cancel, mathtools). The delivered numbering is the unit's own. Assets: no retained span contains a figure environment, a graphics-inclusion command, a PDF-page inclusion, a table or a TikZ picture; no image file is copied into this package, the package declares no asset dependency and no image file is redistributed with it. Licence: version arXiv:2507.20477v1 of this article is distributed under the Creative Commons Attribution 4.0 International licence, as recorded in the arXiv licence metadata for this exact version and shown on the abstract page https://arxiv.org/abs/2507.20477v1. A full rights notice is delivered at licenses/arxiv-2507.20477-CC-BY-4.0.txt. Characters: every retained excerpt is pure ASCII, so the delivered engine, XeLaTeX with its default Latin Modern text font, reports no missing character.
\begin{lemma}\label{lemma: received latent feature}
    Given a set of mapping functions $\{\mathcal{C}_k(\cdot)\}$ that satisfies $\mathcal{C}_k^{-1}(\mathcal{C}_m(\mathbf{f}))\sim\mathcal{N}(0,\mathbf{I})$, and $\mathcal{C}_k^{-1}(\mathcal{C}_k(\mathbf{f}))=\mathbf{f}$,  the received latent feature at user $k$ can be expressed as follows: 
    %For a mapping function that satisfies $\mathcal{C}_k^{-1}(\mathcal{C}_m(\mathbf{f}))\sim\mathcal{N}(0,\mathbf{I})$, and $\mathcal{C}_k^{-1}(\mathcal{C}_k(\mathbf{f}))=\mathbf{f}$, then the received latent feature at user $k$ can be expressed as follows. 
    \begin{align}\label{eq: expression of received latent feature}
    \hat{\mathbf{f}}_k=\alpha_k\mathbf{f}_k+\tau_k\mathbf{n}_r,
    \end{align}
    where $\alpha_k=|\mathbf{h}_k^H\mathbf{v}_k|$ denotes the target signal power, and $\tau_k=\sqrt{\sum_{m=1,m\neq k}|\mathbf{h}_k^H\mathbf{v}_m|^2+\sigma^2}$ denotes the aggregated interference and  noise power. 
    %where $\alpha_k=|\mathbf{h}_k^H\mathbf{v}_k|$ denotes the signal power, and $\beta_k=\sqrt{\sum_{m=1,m\neq k}|\mathbf{h}_k^H\mathbf{v}_m|^2+\sigma^2}$ denotes the noise power.     
\end{lemma}

\begin{proof}
    First, observe that $\mathcal{C}_k^{-1}$ is linear, i.e., for any vectors $\mathbf{a}$ and $\mathbf{b}$, and any scalar $x\in \mathbb{R}$,  the following properties hold:
    $\mathcal{C}_k^{-1}(\mathbf{a}+\mathbf{b})=\mathcal{C}_k^{-1}(\mathbf{a})+\mathcal{C}_k^{-1}(\mathbf{b})$, and $\mathcal{C}_k(x\mathbf{a})=x\mathcal{C}_k(\mathbf{a})$. 
   
   Given these properties, the received latent feature at user $k$ can be decomposed as follows: 
    %\vspace{-3mm}
   \begin{align}
   \hat{\mathbf{f}}_k&=\underbrace{\mathcal{C}_k^{-1}(e^{-j\phi_{k,k}}\mathbf{h}_k^H\mathbf{v}_k\mathcal{C}_k(\mathbf{f}_k))}_{\mathbf{f}_{k,target}}\notag\\&+\underbrace{\sum_{m=1,m\neq k}^K \mathcal{C}_k^{-1}(e^{-j\phi_{k,k}}\mathbf{h}_k^H\mathbf{v}_m\mathcal{C}_m(\mathbf{f}_m))}_{\mathbf{f}_{k,int}}+\underbrace{\mathcal{C}_k^{-1}(\sigma e^{-j\phi} \mathbf{n}_{r})}_{\mathbf{f}_{k,noise}}.
   \end{align}
   %\vspace{-5mm}
   For the desired signal component $\mathbf{f}_{k,target}$, noting that $\mathbf{h}_k^H\mathbf{v}_k=|\mathbf{h}_k^H\mathbf{v}_k|e^{j\phi_{k,k}}$, we have 
   \begin{align}
   \mathbf{f}_{k,target}&=\mathcal{C}_k^{-1}(|\mathbf{h}_k^H\mathbf{v}_k|\mathcal{C}_k(\mathbf{f}_k))=|\mathbf{h}_k^H\mathbf{v}_k|\mathcal{C}_k^{-1}(\mathcal{C}_k(\mathbf{f}_k))=\alpha_k\mathbf{f}_k.
   \end{align}
   For the interference component $\mathbf{f}_{k,int}$, denote $\mathbf{h}_k^H\mathbf{v}_m$ as $\mathbf{h}_k^H\mathbf{v}_m=|\mathbf{h}_k^H\mathbf{v}_m|e^{j\phi_{k,m}}$, then we have 
   \vspace{-3mm}
   \begin{align}
   \mathbf{f}_{k,int}&=\sum_{m=1,m\neq k}^K\mathcal{C}_k^{-1}(e^{j(\phi_{k,m}-\phi_{k,k})}|\mathbf{h}_k^H\mathbf{v}_m|\mathcal{C}_m(\mathbf{f}_m))\\&=\sum_{m=1,m\neq k}^K |\mathbf{h}_k^H\mathbf{v}_m|\mathcal{C}_k^{-1}(e^{j(\phi_{k,m}-\phi_{k,k})}\mathcal{C}_m(\mathbf{f}_m))\\
   &=\sum_{m=1,m\neq k}^K{\bigg[}\underbrace{\cos(\phi_{k,m}-\phi_{k,k})|\mathbf{h}_k^H\mathbf{v}_m|\mathcal{C}_{k}^{-1}(\mathcal{C}_m(\mathbf{f}_m))}_{(a)}\notag\\&+\underbrace{\sin(\phi_{k,m}-\phi_{k,k})|\mathbf{h}_k^H\mathbf{v}_m|\mathcal{C}_{k}^{-1}j(\mathcal{C}_m(\mathbf{f}_m))}_{(b)}{\bigg]}
   \end{align}
   Given that $\mathcal{C}_k^{-1}(\mathcal{C}_m(\mathbf{f}_m))\sim\mathcal{N}(0,\mathbf{I})$, the first term $(a)$ follows the distribution of $\mathcal{N}(0,\cos^2(\phi_{k,m}-\phi_{k,k})|\mathbf{h}_k^H\mathbf{v}_m|^2\mathbf{I})$. For the second term, according to the definition of $\mathcal{C}_k$ and $\mathcal{C}_k^{-1}$, we have $\mathcal{C}_k^{-1}(j\mathcal{C}_m(\mathbf{f}_m))=\mathbf{P}\mathcal{C}_k^{-1}(\mathcal{C}_m(\mathbf{f}_m))$, 
   where $\mathbf{P}=\begin{bmatrix} \mathbf{0}&-\mathbf{I}_{N}\\\mathbf{I}_N&\mathbf{0} \end{bmatrix}$. 
   Therefore, we have $\mathcal{C}_k^{-1}(j\mathcal{C}_m(\mathbf{f}_m))\sim \mathcal{N}(\mathbf{P}\mathbf{0},\mathbf{P}\mathbf{I}\mathbf{P}^T)=\mathcal{N}(\mathbf{0},\mathbf{I})$, indicating that the second term $(b)$ also contributes a Gaussian component  that follows the distribution of $\mathcal{N}(0,\sin^2(\phi_{k,m}-\phi_{k,k})|\mathbf{h}_k^H\mathbf{v}_m|^2\mathbf{I})$. 
   
   By the independence and additivity of Gaussian components, the aggregated interference $\mathbf{f}_{k,int}$ is distributed as $\mathcal{N}(\mathbf{0},\mathbf{I}\sum_{m=1,m\neq k}^K|\mathbf{h}_k^H\mathbf{v}_m|^2)$. 
   Similarly, we have $\mathbf{f}_{k,noise}\sim \mathcal{N}(0,\sigma^2\mathbf{I})$ due to the invariance of the mapping to Gaussian noise. 
   
   Combining the interference and noise, we can conclude that the overall received latent feature at user $k$ can be succinctly expressed as (\ref{eq: expression of received latent feature}), which ends the proof. %\hfill$\Box$
\end{proof}

\end{document}
